% =====================================================================
%  1. INTRODUCTION
% =====================================================================
\section{Introduction}\label{sec:intro}

Le développement logiciel est entré dans l’ère \emph{agent-first} : une part
croissante du code est écrite, relue et modifiée par des modèles de langage
plutôt que par des humains. Les études empiriques convergent sur l’ampleur du
phénomène (gains de productivité mesurés ou déclarés chez les utilisateurs
d’assistants de complétion \cite{ziegler2022,peng2023}, adoption massive et
réorganisation des pratiques de lecture et de vérification du code généré
\cite{vaithilingam2022,barke2023}) tandis que les revues systématiques
recensent déjà plusieurs centaines de travaux sur les modèles de langage
appliqués au génie logiciel \cite{hou2024,fan2023}. Cette bascule remet en
question une hypothèse implicite de nos outils : que le code source est optimisé
pour la \emph{lecture humaine}. Compilateurs, systèmes de construction et
langages ont été conçus autour de cette hypothèse ; si le lecteur principal
devient une machine, l’empilement mérite un réexamen systématique plutôt qu’une
révision marginale.

Or la question qui se pose au praticien n’est pas celle que mesurent la
plupart de ces travaux. La littérature évalue des \emph{capacités} (le modèle
peut-il produire un code correct pour une tâche donnée~?
\cite{chen2021,austin2021}) ou des \emph{effets} (la productivité change-t-elle~?
\cite{peng2023}). Elle documente aussi les modes d’échec : code plausible mais
inexact, dépendances hallucinées, conformité superficielle à la demande
\cite{hou2024}. Mais entre le modèle et le produit livré s’étend un territoire
peu exploré~: le \emph{processus} — l’ensemble des règles qui déterminent ce
qu’un agent a le droit d’écrire, quand une écriture devient permanente, et
comment un projet survit au changement d’agent. C’est ce territoire que
TRIADE se propose de gouverner.

TRIADE (\emph{Intention, Binaire, Carte}) est un protocole de gouvernance du
développement logiciel conçu pour un monde où l’exécutant est un agent. Son
principe est simple à énoncer~: \emph{aucune écriture n’est permanente sans
vérification, et aucune évolution ne commence sans intention écrite}. Toute
modification traverse trois artefacts versionnés — la couche I
(spécifications d’intention diffables), la couche B (l’artefact exécutable
scellé), la couche C (une carte indexée, signée, qui relie les deux) — et dix
règles à sanction mécanique (refus de construction, refus de commit, blocage
de session) plutôt qu’à recommandation. Un verrou croisé cryptographique relie
la carte et le binaire dans les deux sens ; une fonction objectif
lexicographique (sécurité $>$ performance $>$ optimisation $>$ coût en tokens)
arbitre tout conflit sans compensation possible entre rangs. Le protocole
existe en version~0.3, spécifié dans le dépôt du projet avec son journal de
décisions.

Cet article restructure cette spécification selon les standards d’un article de
recherche en génie logiciel et la confronte à sa première validation à
l’échelle~: le recodage complet de RoomEq, un égaliseur acoustique temps réel
d’environ 15~kloc C\#/.NET~8, exécuté en sept vagues depuis une base scellée
v2.2.2. Nos contributions sont les suivantes~:
\begin{itemize}
\item \textbf{C1 — Formalisation.} Nous donnons au protocole une présentation
formalisée~: état du projet comme quadruplet, prédicat de cohérence du verrou
croisé, règles de transformation en système de transitions, oracle
différentiel, fonction objectif lexicographique
(section~\ref{sec:methode}).
\item \textbf{C2 — Étude de cas à l’échelle.} Nous rapportons le recodage
complet d’une application réelle sous gouvernance TRIADE, avec ses métriques
(structure, défense, gouvernance), ses comparatifs de code avant/après et sa
séquence de publication — y compris deux échecs d’intégration continue honorés
puis corrigés (section~\ref{sec:etude}).
\item \textbf{C3 — Reverse-specs.} Nous introduisons la technique de
spécification rétro-active de binaires scellés hérités, qui étend la couche I
au patrimoine existant sans le recoder (section~\ref{sec:reverse}).
\item \textbf{C4 — Auto-évaluation critique.} Nous délimitons explicitement ce
que la démonstration établit et ne prouve pas, selon le cadre des menaces sur
la validité de Wohlin et al. \cite{wohlin2012} — notamment la validité externe
d’un cas unique et l’écart entre refactorisation source et génération binaire
native (section~\ref{sec:menaces}).
\end{itemize}

L’article suit l’ordre attendu~: la section~\ref{sec:connexes} situe TRIADE
par rapport à la dette technique, à la conception par contrat et au génie
logiciel assisté par IA~; la section~\ref{sec:methode} formalise le
protocole~; la section~\ref{sec:etude} présente l’étude de cas~; les
sections~\ref{sec:discussion} et~\ref{sec:menaces} discutent les résultats et
leurs limites~; la section~\ref{sec:conclusion} conclut. Le lecteur pressé
trouvera en tableau~\ref{tab:regles} la traduction de chaque règle dans le
dépôt de l’étude de cas — c’est la table qui a gouverné sept vagues de
recodage.
